% Accessible LaTeX Math Test Template
% Purpose: Diagnostically rich math expressions for testing screen reader
%          speech output. Expressions are chosen so that character-by-character
%          reading sounds obviously wrong, making it easy for an untrained
%          tester to identify incorrect AT behavior.
%
% Compiler: LuaLaTeX (required for unicode-math and PDF tagging)
% Tested with: Rolling TeX Live (Labs), Overleaf
%
% -----------------------------------------------------------------------
% \DocumentMetadata MUST come before \documentclass
% -----------------------------------------------------------------------
\DocumentMetadata{
  lang=en,
  tagging=on,
  pdfstandard = {ua-2, A-4f},
  tagging-setup={
    math/setup = {mathml-AF, mathml-SE}
  }
}

\documentclass{article}
\usepackage[margin=1in]{geometry}
\usepackage{parskip} 

% graphicx required for \includegraphics with alt= key
\usepackage{graphicx}

% unicode-math replaces amsmath and amssymb.
% Must be loaded AFTER graphicx.
\usepackage{unicode-math}

\usepackage{etoolbox}

\makeatletter
\AtBeginDocument{%
  \patchcmd{\@maketitle}
    {\null\vskip 2em}
    {}%
    {}%
    {\PackageWarningNoLine{title-spacing}
      {Could not remove the article title top skip from \string\@maketitle}}%
}
\makeatother

\title{Tagged PDF with MathML Test: AF+SE}
\author{Brian Richwine, IU}
\date{}

% -----------------------------------------------------------------------
% Semantic macro for limits: makes the operand structurally explicit
% in the LaTeX source and in MathML output.
% Usage: \limitof{var}{point}{expression}
% -----------------------------------------------------------------------
\newcommand{\limitof}[3]{\lim_{#1 \to #2}\!\left(#3\right)}

\begin{document}
\raggedright

\maketitle

% -----------------------------------------------------------------------
\section*{Template notes}
% -----------------------------------------------------------------------

\textbf{Configuration:} MathML embedded directly in the PDF tag tree as Structure Elements

% -----------------------------------------------------------------------
% PURPOSE OF THIS DOCUMENT
% -----------------------------------------------------------------------

This document contains math expressions chosen to make incorrect screen 
reader behavior immediately obvious, even to a tester who is not familiar 
with math accessibility, math speech grammars, etc. 

For each expression, a sample expected speech output is shown alongside 
a sample the incorrect output. The expected speech text is modeled after 
the speech produced by the JAWS screen reader version 2026's default 
settings when it is rendering MathML as speech text. The incorrect speech 
text is modeled after the JAWS screen reader when it reads only the visible 
rendered characters with no structural information from MathML. 

If the screen reader output more closely matches the sample ``incorrect'' 
speech, the math encoding is not being handled correctly somewhere in 
the PDF viewer, OS, and assistive technology stack. 

\textbf{Author's note on the sample speech text:} The sample speech texts 
are from JAWS 2026 with the default settings. The are not provided as 
exemplars of how the expressions should be spoken. I would not recommend 
the default settings be used. Instead I highly recommend switching to 
either ClearSpeak Verbose until the user has mastered depending upon speech 
text for math access. 

% -----------------------------------------------------------------------
\section*{1. Fractions}
% -----------------------------------------------------------------------

\subsection*{1a. Simple fraction (inline)}

The slope of a line is given by 
$m = \frac{\Delta y}{\Delta x}$.

\textbf{Expected speech:} ``m is equal to the fraction with numerator 
cap delta y and denominator cap delta x end fraction MathContent'' 
(or similar phrasing depending on screen reader)

\textbf{Incorrect speech (MathML not working):} ``italics small m 
equals delta italics small y delta italics small x''

In the incorrect speech example the fraction is not announced; the 
numerator and denominator are read as a flat sequence with nothing 
indicating one is over the other.

\subsection*{1b. Quadratic formula (display)}

The solutions to $ax^2 + bx + c = 0$ are given by the quadratic formula:
\[
  x = \frac{-b \pm \sqrt{b^2 - 4ac}}{2a}
\]

\textbf{Expected speech:} ``x equals fraction, negative b plus or minus 
square root, b squared minus 4 eigh c, end root over, 2 eigh, 
end fraction MathContent''

\textbf{Incorrect speech (MathML not working):} ``italics small x equals 
minus italics small b plus or minus italics small b 2 minus 4 italics 
small a italics small c 2 italics small a''

In the incorrect speech example the fraction and the square root grouping 
are not announced. The numerator and denominator merge into a single flat 
string, and the radical over $b^2 - 4ac$ is not mentioned.

\subsection*{1c. Nested fraction (display)}

\[
  \frac{1}{1 + \dfrac{1}{1 + \dfrac{1}{x}}}
\]

\textbf{Expected speech:} ``fraction, 1 over, 1 plus fraction, 1 over, 1 plus, 
fraction, 1 over x, end fraction end fraction end fraction MathContent''

\textbf{Incorrect speech (MathML not working):} ``1 1 plus 1 1 plus 1 italics 
small x'' 

In the incorrect speech example every fraction is not announced. All six 
numerals and the variable collapse into a flat sequence with no indication 
of nesting or structure whatsoever.

% -----------------------------------------------------------------------
\section*{2. Square Roots and Radicals}
% -----------------------------------------------------------------------

\subsection*{2a. Simple square root (inline)}

The distance between two points is $d = \sqrt{(x_2 - x_1)^2 + (y_2 - y_1)^2}$.

\textbf{Expected speech:} ``d equals square root, open x 2, 
minus x 1 close squared plus, open y 2, minus y 1 close squared, 
end root MathContent''

\textbf{Incorrect speech (MathML not working):} ``italics small d equals 
left paren italics small x 2 minus italics small x 1 right paren 
2 plus left paren italics small y 2 minus italics small y 1 right paren 2''

In the incorrect speech example the square root symbol is silent; 
the radicand is read as if it were ordinary text following the equals sign, 
with no indication that the entire expression is under a radical. The
exponents are not announced as such.

\subsection*{2b. Fraction under a radical (display)}

\[
  \sqrt{\frac{x + 1}{x - 1}}
\]

\textbf{Expected speech:} ``square root, fraction, x plus 1, over, 
x minus 1, end fraction end root MathContent''

\textbf{Incorrect speech (MathML not working):} ``italics small x plus 1 
italics small x minus 1''

Both the radical and the fraction bar are silent. The rendered glyphs are
just the four terms and the two operator signs; nothing indicates they form
a fraction nested inside a square root.

\subsection*{2c. nth root (display)}

\[
  \sqrt[n]{a^n} = \lvert a \rvert \quad \text{for even } n
\]

\textbf{Expected speech:} ``n-th root eigh to the n-th, end root equals, 
absolute value eigh, for even n MathContent''

\textbf{Incorrect speech (MathML not working):} ``italics small a italics small n 
italics small n equals vertical bar italics small a vertical bar for even italics small n''

The radical index, the radical itself, and the exponent are not announced as such. 

% -----------------------------------------------------------------------
\section*{3. Exponents and Subscripts}
% -----------------------------------------------------------------------

\subsection*{3a. Nested exponents (inline)}

Consider the tower $2^{2^n}$.

\textbf{Expected speech:} ``2 raised to the 2 to the n-th power MathContent''

\textbf{Incorrect speech (MathML not working):} ``2 2 italics small n''

All three characters are read as a flat sequence. The two levels of
exponentiation are completely invisible; a listener hearing ``2 2 italics small n'' has
no way to recover the intended meaning.

\subsection*{3b. Exponential function with compound exponent (display)}

\[
  f(x) = e^{-x^2 / 2}
\]

\textbf{Expected speech:} ``f of x equals, e raised to the negative x squared, 
divided by 2 power MathContent''

\textbf{Incorrect speech (MathML not working):} ``italics small f left paren 
italics small x right paren italics small e minus italics small x 2 slash 2''

The function speech for $f(x)$, the exponent structure, and
the fraction in the exponent are not announced as such. The rendered child
characters---f, x, e, x, 2, 2---are read as a meaningless flat string.

\subsection*{3c. Subscripts and superscripts together (display)}

\[
  \sum_{i=1}^{n} i^2 = \frac{n(n+1)(2n+1)}{6}
\]

\textbf{Expected speech:} ``sum from i equals 1 to n of, i squared equals fraction, 
n times, open n plus 1 close, times, open 2 n plus 1, close, over 6, end fraction  
MathContent''

\textbf{Incorrect speech (MathML not working):} ``sigma italics small i 1 italics 
small n italics small i 2 equals italics small n left paren italics small n plus 
1 right paren left paren 2 n plus 1 right paren 6''

The summation's lower limit ($i=1$), upper limit ($n$),
and summand ($i^2$) merge with the right-hand side into an undifferentiated
run of characters. The fraction bar on the right-hand side is also silent.

% -----------------------------------------------------------------------
\section*{4. Limits and Integrals}
% -----------------------------------------------------------------------

\subsection*{4a. Limit with explicit operand (display)}

\[
  \limitof{x}{0}{\frac{\sin x}{x}} = 1
\]

\textbf{Expected speech:} ``limit as x approaches 0 open, fraction, sine x, 
over x, end fraction close equals 1 MathContent''

\textbf{Incorrect speech (MathML not working):} ``lim italics small x 
rightwards arrow 0 left paren sin right paren italics small x 
italics small x 1''

The arrow ($\to$) may be spoken as a symbol name or silently
dropped. The fraction ``$\sin x$ over $x$'' collapses to ``sin x x'' with
no indication that the sine expression is in the numerator and $x$ alone
is the denominator.

\subsection*{4b. Definite integral (display)}

\[
  \int_{0}^{1} x^2 \, dx = \frac{1}{3}
\]

\textbf{Expected speech:} ``integral from 0 to 1 of, x squared d x equals 1 third  
MathContent''

\textbf{Incorrect speech (MathML not working):} ``integral blank 0 1 italics small x 
2 italics small d italics small x 1 3''

The limits of integration (0 and 1) are
read as bare numbers with no indication they are bounds. The fraction
``$\frac{1}{3}$'' collapses to ``1 3'' with no spoken relationship between
the two digits.

% -----------------------------------------------------------------------
\section*{5. Matrix}
% -----------------------------------------------------------------------

\subsection*{5a. 2 × 2 matrix with non-trivial entries (display)}

\[
  A = \begin{pmatrix} a & b \\ c & d \end{pmatrix}
\]

\textbf{Expected speech:} ``cap eigh equals the 2 by 2 matrix row 1 eigh, b row 2 c, d  
MathContent'' (exact phrasing varies by screen reader and math speech
engine).

\textbf{Incorrect speech (MathML not working):} ``italics capital A equals left paren italics small a italics small b italics small c italics small d right paren''.
The matrix delimiters and the row/column structure are entirely absent.
The four entries are read as a flat list with no indication they form a
rectangular array, let alone that $a$ and $b$ are in the same row while
$c$ and $d$ are in another.

\subsection*{5b. Matrix equation showing determinant (display)}

\[
  \det(A) = \begin{vmatrix} a & b \\ c & d \end{vmatrix} = ad - bc
\]

\textbf{Expected speech:} ``det of cap eigh, equals, the 2 by 2 determinant; 
row 1; eigh, b; row 2; c, d; equals, eigh d minus b c MathContent''

\textbf{Incorrect speech (MathML not working):} ``det left paren italics capital A right paren equals vertical bar italics small a italics small b italics small c italics small d vertical bar equals italics small a italics small d minus italics small b italics small c''

The vertical bars are silent. The entries are read as a flat continuation
of the surrounding expression, indistinguishable from the ``$ad - bc$''
that follows. A tester hearing ``a b c d a d b c'' cannot tell where the
matrix ends and the scalar expression begins.

% -----------------------------------------------------------------------
\section*{6. Combined Structure}
% -----------------------------------------------------------------------

\subsection*{6a. Euler's identity (display)}

\[
  e^{i\pi} + 1 = 0
\]

\textbf{Expected speech:} ``e raised to the i pi power plus 1 
equals 0 MathContent''

\textbf{Incorrect speech (MathML not working):} ``italics small e 
italics small i italics small pi plus 1 equals 0''

The exponent structure is silent; $i\pi$ is read as if it were at the
same level as the $+\,1$ that follows rather than as the exponent of $e$.

\subsection*{6b. Normal distribution probability density (display)}

\[
  f(x) = \frac{1}{\sigma\sqrt{2\pi}}\, e^{-\frac{(x-\mu)^2}{2\sigma^2}}
\]

\textbf{Expected speech:} ``f of x equals fraction, 1 over, sigma times, 
square root, 2 pi, end root end fraction times e raised to the negative fraction, 
open x minus mu close squared, over, 2 sigma squared, end fraction power  
MathContent''

\textbf{Incorrect speech (MathML not working):} ``italics small f left paren 
italics small x right paren equals 1 italics small sigma 2 italics small pi 
italics small e minus left paren x minus italics small mu right paren 2 2 
italics small sigma 2''

Every structural element is silent: the outer fraction bar, the radical,
the exponent grouping, and the inner fraction in the exponent all collapse.

\end{document}